\begin{resumo}

A análise computacional de microbiomas requer a integração de múltiplas
ferramentas estatísticas, formatos de dados heterogêneos e pipelines
bioinformáticos complexos — um cenário que dificulta o trabalho de
pesquisadores sem formação em programação. Este trabalho descreve o
projeto e a implementação de uma plataforma web acadêmica voltada à
análise integrada de micobioma e metataxonômica, desenvolvida para
suportar três projetos científicos paralelos: \textbf{INOVAHERB}
(análise fatorial de micobioma ITS com ANCOM-BC2), \textbf{Pós-Fogo}
(série temporal de recuperação de solo via 16S com MaAsLin2) e
\textbf{Biorremediação} (correlação de microbiota com variáveis
edáficas). A arquitetura adotada baseia-se em microsserviços: um
backend FastAPI (Python) expondo uma API REST com padrão CQRS, um R
Worker assíncrono que executa as análises estatísticas pesadas
(DESeq2, ANCOM-BC2, MaAsLin2, SpiecEasi, Random Forest, GSEA,
FUNGuild e PICRUSt2), um frontend Next.js 14 com visualizações
interativas via Plotly.js e Cytoscape.js, e uma camada de dados
composta por PostgreSQL 16 (incluindo Large Objects para armazenamento binário) e Elasticsearch 8. A orquestração
de jobs é realizada por meio de filas nativas do PostgreSQL
(\texttt{LISTEN/NOTIFY} com \texttt{FOR UPDATE SKIP LOCKED}), sem
dependência de sistemas de mensageria externos. A plataforma é
implantada em cluster Kubernetes leve (k3s) com cinco nós AMD e
suporta autenticação institucional via Google OAuth restrita ao domínio
\texttt{@ufvjm.edu.br}. Os resultados esperados incluem a entrega de um
painel unificado com seis PCoAs, duas por projeto, gerado como figura
central do TCC após a conclusão dos três pipelines.

\vspace{\onelineskip}
\noindent
\textbf{Palavras-chave}: bioinformática; micobioma; metataxonômica; plataforma web; FastAPI; R Worker; PostgreSQL; QIIME2.

\end{resumo}
